%% 10-intro.tex — Introduction

\section{Introduction}
\label{sec:introduction}

The German Digital Library (Deutsche Digitale Bibliothek, DDB)\footnote{DDB, \url{https://www.deutsche-digitale-bibliothek.de}} aggregates over 65M objects\footnote{of which 27M are digitized} from German memory institutions --- archives, libraries, museums, media libraries, and monument preservation offices, referred to hereafter as \emph{sectors} --- all described using the Europeana Data Model (EDM)\footnote{EDM, \url{https://pro.europeana.eu/page/edm-documentation}}.
The corpus has no public SPARQL endpoint, no Linked Data browser, and no machine-readable interface for KG-based research.

Converting DDB-EDM records to a queryable knowledge graph (KG) is complicated by sector-specific metadata conventions that diverge across institutional traditions.
Each sector applies locally-specific field interpretations and carries cataloging practices into every record; no cross-sector authority governs how EDM properties are instantiated.
This variation is invisible to schema-level or structure-based alignment tools; corpus-based analysis is therefore a precondition for reliable automated alignment.

%GEMEA (German Memory Atlas) is a publicly accessible KG over 27 million digitized objects from these institutions, and the first open, SPARQL-queryable structured access layer over the DDB corpus.
Three contributions are described in this paper:

\begin{enumerate}
    \item \textbf{Resource.} GEMEA /\textipa{g@"meI.j@}/
(\textbf{Ge}rman \textbf{Me}mory \textbf{A}tlas), a publicly queryable KG of 27M German cultural heritage objects, served via a QLever SPARQL endpoint~\cite{bast2017qlever} and browsable through a SHMARQL~\cite{shmarql} Linked Data interface\footnote{\url{https://github.com/epoz/shmarql}} --- the first open, SPARQL-accessible KG over the DDB corpus.

    \item \textbf{Ontology alignment Layer.} A \textbf{MOCHO} \footnote{Middle Ontology for Cultural Heritage Objects-}-based\cite{tan2026navigating} /\textipa{"mo.tSo}/ alignment layer mapping DDB-EDM metadata to RDA/FRBR-compatible representations, demonstrated on a representative subset, with alignment gaps documented as findings rather than omitted.

    \item \textbf{Agentic access layer.} An MCP/MCPO interface on the QLever endpoint, enabling LLM-based agents to issue SPARQL queries, retrieve entity descriptions, and integrate GEMEA into multi-step reasoning workflows.
\end{enumerate}

As a KG derived from in-use institutional data, GEMEA is intended as a research testbed; Section~\ref{sec:usage} details supported use cases and downstream research scenarios.
Data is licensed CC~BY~4.0; the SPARQL endpoint, Linked Data browser, and source code are available at \url{https://gemea.ise.fiz-karlsruhe.de}.

\paragraph{Paper structure.}
Sections~\ref{sec:related}--\ref{sec:sustainability} cover related work, resource description, quality, usage, impact, and sustainability in turn.
